%===============================================================================
\chapter{MAT1093 --- Precalculus}
\label{chap:unit-2531}
%===============================================================================

\ChapterWarning

%===============================================================================
\section{Section 1.1 - Functions and Function Notation}
\label{sec:sect-c5c2}
%===============================================================================

We summarise the core toolkit (library) functions every student of college
algebra and trigonometry should know by name, formula, and graph.  For each
function the domain and range are stated in interval notation.

\subsection*{Constant Function}

For the constant function $f(x) = c$ the domain consists of all real numbers;
there are no restrictions on the input.  The only output value is the constant
$c$, so the range is the set $\{c\}$.  In interval notation this is written
$[c,c]$, the interval that both begins and ends with $c$.

\[
  \text{Domain: }(-\infty, \infty) \qquad \text{Range: }[c, c]
\]

\subsection*{Identity Function}

For the identity function $f(x) = x$ there is no restriction on $x$.  Both
the domain and range are the set of all real numbers.

\[
  \text{Domain: }(-\infty, \infty) \qquad \text{Range: }(-\infty, \infty)
\]

\subsection*{Absolute Value Function}

For the absolute value function $f(x) = |x|$ there is no restriction on $x$.
The outputs are $x$ for $x \geq 0$ and $-x$ for $x < 0$, so the range is all
numbers greater than or equal to zero.

\[
  \text{Domain: }(-\infty, \infty) \qquad \text{Range: }[0, \infty)
\]

\subsection*{Quadratic Function}

For the quadratic function $f(x) = x^{2}$ the domain is all real numbers.
Because the graph does not include negative $y$-values, the range contains
only nonnegative real numbers.

\[
  \text{Domain: }(-\infty, \infty) \qquad \text{Range: }[0, \infty)
\]

\subsection*{Cubic Function}

For the cubic function $f(x) = x^{3}$ both the domain and range include all
real numbers, since the graph extends without bound in both the horizontal and
vertical directions.

\[
  \text{Domain: }(-\infty, \infty) \qquad \text{Range: }(-\infty, \infty)
\]

\subsection*{Rational (Reciprocal) Function}

For $f(x) = \frac{1}{x}$ we cannot divide by zero, so $x = 0$ is excluded
from the domain.  Moreover, $\frac{1}{x}$ can never equal zero, so $0$ is
also excluded from the range.

\[
  \text{Domain: }(-\infty, 0) \cup (0, \infty) \qquad
  \text{Range: }(-\infty, 0) \cup (0, \infty)
\]

\subsection*{Squared Reciprocal Function}

For $f(x) = \frac{1}{x^{2}}$ the same exclusion of $x = 0$ applies.
Additionally, because $x^{2} > 0$ for all $x \neq 0$, the outputs are always
positive, so the range contains only positive real numbers.

\[
  \text{Domain: }(-\infty, 0) \cup (0, \infty) \qquad
  \text{Range: }(0, \infty)
\]

\subsection*{Square Root Function}

For $f(x) = \sqrt{x}$ we cannot take the square root of a negative real
number, so the domain is restricted to $x \geq 0$.  Because the principal
square root is non-negative, the range is also $[0, \infty)$.

\[
  \text{Domain: }[0, \infty) \qquad \text{Range: }[0, \infty)
\]

\subsection*{Cube Root Function}

For $f(x) = \sqrt[3]{x}$ there is no restriction: cube roots of negative
numbers are real and negative.  Both the domain and range are all real
numbers, and the function is odd.

\[
  \text{Domain: }(-\infty, \infty) \qquad \text{Range: }(-\infty, \infty)
\]

\begin{itemize}
  \item \href{https://mathresearch.utsa.edu/wikiFiles/MAT1053/Domain_Range_and_Toolkit_Functions/MAT1053_M1.2Domain_Range_and_Toolkit_Functions.pdf}{Domain, Range, and Toolkit Functions}, UTSA book chapter
  \item \href{https://courses.lumenlearning.com/cuny-hunter-collegealgebra/chapter/toolkit-functions/}{Toolkit Functions}, Lumen Learning
\end{itemize}

%-------------------------------------------------------------------------------
\subsection{Learning Objectives}
\label{ssec:lear-fe6e}
%-------------------------------------------------------------------------------

\begin{enumerate}
    \item Determine whether a relation represents a function.
    \item Find the value of a function.
    \item Determine whether a function is one-to-one.
    \item Use the vertical line test to identify functions.
    \item Graph the functions listed in the library of functions.
\end{enumerate}

%-------------------------------------------------------------------------------
\subsection{Assessment Instrument}
\label{ssec:asse-3669}
%-------------------------------------------------------------------------------

Assess mastery of toolkit functions by asking students to identify the domain and range of
each library function, sketch its graph from memory, and explain the key feature that
distinguishes one toolkit function from another (e.g., why the square root function has a
restricted domain while the cube root does not).

\textbf{Example problems.}
\begin{enumerate}
    \item State the domain and range of each toolkit function: $f(x) = c$, $f(x) = |x|$, $f(x) = x^{2}$, $f(x) = \sqrt{x}$, and $f(x) = \frac{1}{x}$.
    \item Sketch the graphs of $f(x) = x^{3}$ and $g(x) = \sqrt[3]{x}$ on the same axes, state their domains and ranges, and describe the symmetry of each graph.
    \item A student claims that $f(x) = \frac{1}{x^{2}}$ has the same range as $f(x) = \frac{1}{x}$. Determine whether this claim is correct and justify your answer.
\end{enumerate}

%===============================================================================
\section{Section 1.2 - Domain and Range}
\label{sec:sect-5062}
%===============================================================================

A function $f$ is a mapping from a domain set $A$ to a range set $B$ such that
every element $a \in A$ is paired with exactly one element $b \in B$, written
$f : A \to B$.  The set of all permissible inputs is the \textbf{domain}; the
set of all resulting outputs is the \textbf{range}.

\subsection*{One-to-One Functions and the Horizontal Line Test}

A function is \textbf{one-to-one} (injective) if distinct inputs always
produce distinct outputs.  Geometrically, no horizontal line may intersect
the graph more than once.  The algebraic test states: $f$ is one-to-one if
$f(a) = f(b)$ implies $a = b$ for all $a, b$ in the domain.

\textbf{Example.} Prove that $f(x) = \frac{1-2x}{1+x}$ is injective.

Assume $f(a) = f(b)$.  Then
\begin{align}
  \frac{1-2a}{1+a} &= \frac{1-2b}{1+b} \\
  (1+b)(1-2a) &= (1+a)(1-2b) \\
  1 - 2a + b - 2ab &= 1 - 2b + a - 2ab \\
  -2a + b &= -2b + a \\
  3b &= 3a \\
  a &= b.
\end{align}
Because $f(a) = f(b)$ forces $a = b$, the function is injective.

By contrast, $g(x) = x^{2}$ is \emph{not} injective: the same argument leads
to $a = \pm b$, so two distinct inputs can share the same output.

\begin{itemize}
  \item \href{https://en.wikibooks.org/wiki/Calculus/Functions}{Functions, Wikibooks: Calculus}
  \item \href{https://www.math.uh.edu/~jiwenhe/Math1432/lectures/lecture01_handout.pdf}{One-to-One Functions and Inverses, University of Houston}
\end{itemize}

%-------------------------------------------------------------------------------
\subsection{Learning Objectives}
\label{ssec:lear-fe6e}
%-------------------------------------------------------------------------------

\begin{enumerate}
    \item Find the domain of a function defined by an equation.
    \item Graph piecewise-defined functions.
\end{enumerate}

%-------------------------------------------------------------------------------
\subsection{Assessment Instrument}
\label{ssec:asse-3669}
%-------------------------------------------------------------------------------

Evaluate understanding of domain, range, and one-to-one functions through problems that
require students to find the domain of expressions with algebraic restrictions, apply the
horizontal line test, and prove or disprove injectivity using the algebraic definition.

\textbf{Example problems.}
\begin{enumerate}
    \item Find the domain of $f(x) = \frac{1-2x}{1+x}$ and determine whether $f$ is a one-to-one function using both the horizontal line test and the algebraic injectivity test.
    \item Prove that $f(x) = 3x - 7$ is injective by showing that $f(a) = f(b)$ implies $a = b$.
    \item Determine all values of $x$ for which the piecewise function $f(x) = \begin{cases} x^{2} & x \geq 0 \\ -x & x < 0 \end{cases}$ is defined, and state its domain and range.
\end{enumerate}

%===============================================================================
\section{Section 2.1 - Angles}
\label{sec:sect-7b60}
%===============================================================================

A \textbf{sector} is a region bounded by two radii of a circle and the arc
between them.  The angle that determines the sector is called the
\textbf{central angle}.

\subsection*{Degree Measure}

A full rotation is divided into $360$ degrees ($360^{\circ}$).  Each degree
contains $60$ minutes ($'$) and each minute contains $60$ seconds ($''$).  To
convert $X^{\circ}\,Y'\,Z''$ to decimal degrees:
\[
  X + \frac{Y}{60} + \frac{Z}{3600}.
\]
For example, $23^{\circ}\,18'\,38'' = 23.3106^{\circ}$.

Angles are classified by their measure: an \textbf{acute} angle measures
between $0^{\circ}$ and $90^{\circ}$; a \textbf{right} angle measures exactly
$90^{\circ}$; an \textbf{obtuse} angle measures between $90^{\circ}$ and
$180^{\circ}$.

\subsection*{Radian Measure}

A \textbf{radian} is the angle subtended at the centre of a circle by an arc
equal in length to the radius.  Since the full circumference is $2\pi r$,
there are $2\pi$ radians in a complete rotation.  The conversion relations are
\[
  1^{\circ} = \frac{\pi}{180}\ \text{rad}, \qquad
  1\ \text{rad} = \frac{180}{\pi}\ {}^{\circ}.
\]
To convert $20^{\circ}$ to radians: $20^{\circ} = \frac{20\pi}{180} =
\frac{\pi}{9}\ \text{rad}$.

\subsection*{Arc Length}

For a central angle $\theta$ (in radians) and radius $r$, the arc length is
\[
  s = \theta r.
\]
For a sector with $r = 6\ \text{cm}$ and $\theta = 53^{\circ} =
\frac{53\pi}{180}\ \text{rad}$,
\[
  s = \frac{53\pi}{180} \times 6 \approx 5.55\ \text{cm}.
\]

\subsection*{Area of a Sector}

The area of a sector with radius $r$ and central angle $\theta$ (in radians)
is
\[
  A = \frac{1}{2}r^{2}\theta.
\]
For $r = 3\ \text{cm}$ and $\theta = 2\pi$:
$A = \frac{1}{2}(9)(2\pi) \approx 28.3\ \text{cm}^{2}$.

\subsection*{The Unit Circle}

The \textbf{unit circle} is the circle $x^{2} + y^{2} = 1$ centred at the
origin.  Rotating a ray from the positive $x$-axis by angle $\theta$
(counterclockwise for $\theta > 0$) yields a point
$A = (x_{A}, y_{A}) = (\cos\theta, \sin\theta)$ on the circle.  The six
trigonometric functions are then defined for any angle, not just acute angles
in a right triangle.

\begin{itemize}
  \item \href{https://en.wikibooks.org/wiki/Geometry/Angles}{Angles, Wikibooks: Geometry}
  \item \href{https://en.wikibooks.org/wiki/A-level_Mathematics/OCR/C2/Circles_and_Angles}{Circles and Angles, Wikibooks: A-Level Mathematics}
\end{itemize}

%-------------------------------------------------------------------------------
\subsection{Learning Objectives}
\label{ssec:lear-fe6e}
%-------------------------------------------------------------------------------

\begin{enumerate}
    \item Draw angles in standard position.
    \item Convert between degrees and radians.
    \item Find coterminal angles.
    \item Find the length of a circular arc.
    \item Use linear and angular speed to describe motion on a circular path.
\end{enumerate}

%-------------------------------------------------------------------------------
\subsection{Assessment Instrument}
\label{ssec:asse-3669}
%-------------------------------------------------------------------------------

Test fluency with angle measurement by requiring unit conversions between degrees and
radians, conversion of degrees-minutes-seconds notation, and arc-length or sector-area
calculations that require $\theta$ in radians.

\textbf{Example problems.}
\begin{enumerate}
    \item Convert $150^{\circ}$ to radians and convert $\frac{5\pi}{4}$ radians to degrees.
    \item Convert $47^{\circ}\,23'\,15''$ to decimal degrees, then to radians. Round to four decimal places.
    \item A circle has radius $8\,\text{cm}$. Find (a) the arc length and (b) the area of the sector subtended by a central angle of $\frac{2\pi}{3}$ radians.
\end{enumerate}

%===============================================================================
\section{Section 2.2 - The Unit Circle: Sine and Cosine Functions}
\label{sec:sect-008d}
%===============================================================================

Consider a point $(x, y)$ on the unit circle ($r = 1$) at angle $\theta$ in
standard position.  The six trigonometric functions are defined as
\begin{align}
  \cos\theta &= x, & \sec\theta &= \frac{1}{x}, \\
  \sin\theta &= y, & \csc\theta &= \frac{1}{y}, \\
  \tan\theta &= \frac{y}{x}, & \cot\theta &= \frac{x}{y}.
\end{align}
The cosine and sine of any angle are simply the $x$- and $y$-coordinates of
the corresponding point on the unit circle; this extends their definitions
beyond the right-triangle context to include quadrantal angles and angles
greater than $360^{\circ}$.

\subsection*{Values at Key Angles}

The terminal sides of $30^{\circ}$, $45^{\circ}$, and $60^{\circ}$ intersect
the unit circle at the following points:

\[
  \begin{array}{|c|c|c|}
    \hline
    \textbf{Angle} & x\text{-coordinate} & y\text{-coordinate} \\
    \hline
    30^{\circ} & \dfrac{\sqrt{3}}{2} & \dfrac{1}{2} \\[6pt]
    \hline
    45^{\circ} & \dfrac{\sqrt{2}}{2} & \dfrac{\sqrt{2}}{2} \\[6pt]
    \hline
    60^{\circ} & \dfrac{1}{2} & \dfrac{\sqrt{3}}{2} \\[6pt]
    \hline
  \end{array}
\]

For example, $\cos(45^{\circ}) = \frac{\sqrt{2}}{2}$,
$\sin(60^{\circ}) = \frac{\sqrt{3}}{2}$, and
$\tan(30^{\circ}) = \frac{1/2}{\sqrt{3}/2} = \frac{1}{\sqrt{3}}$.

For the quadrantal angle $\theta = 90^{\circ}$, the point is $(0, 1)$, so
$\sin(90^{\circ}) = 1$, $\cos(90^{\circ}) = 0$, and $\cot(180^{\circ})$ is
undefined because the point is $(-1, 0)$ giving $x/y = -1/0$.

\begin{itemize}
  \item \href{https://en.wikibooks.org/wiki/Trigonometry/The_Unit_Circle}{The Unit Circle, Wikibooks: Trigonometry}
  \item \href{https://en.wikibooks.org/wiki/High_School_Trigonometry/Defining_Trigonometric_Functions}{Defining Trigonometric Functions, Wikibooks: High School Trigonometry}
\end{itemize}

%-------------------------------------------------------------------------------
\subsection{Learning Objectives}
\label{ssec:lear-fe6e}
%-------------------------------------------------------------------------------

\begin{enumerate}
    \item Find function values for the sine and cosine of $30^{\circ}$ ($\pi/6$), $45^{\circ}$ ($\pi/4$), and $60^{\circ}$ ($\pi/3$).
    \item Identify the domain and range of sine and cosine functions.
    \item Use reference angles to evaluate trigonometric functions.
\end{enumerate}

%-------------------------------------------------------------------------------
\subsection{Assessment Instrument}
\label{ssec:asse-3669}
%-------------------------------------------------------------------------------

Assess the ability to use the unit circle to evaluate all six trigonometric functions at
standard angles and quadrantal angles, including identifying where functions are undefined
and reading coordinates directly from the circle.

\textbf{Example problems.}
\begin{enumerate}
    \item Use the unit circle to find the exact values of $\sin(150^{\circ})$, $\cos(270^{\circ})$, and $\tan(315^{\circ})$.
    \item Identify the ordered pair on the unit circle corresponding to $\theta = \frac{2\pi}{3}$ and use it to evaluate all six trigonometric functions at that angle.
    \item Explain why $\cot(180^{\circ})$ is undefined, and find all angles $\theta \in [0^{\circ}, 360^{\circ})$ for which $\sec\theta$ is undefined.
\end{enumerate}

%===============================================================================
\section{Section 2.3 - The Other Trigonometric Functions}
\label{sec:sect-d920}
%===============================================================================

Some of the most important trigonometric identities derive directly from the
Pythagorean theorem.  For a right triangle with legs $a$, $b$ and hypotenuse
$c$:
\[
  a^{2} + b^{2} = c^{2}.
\]

\subsection*{Pythagorean Identity}

Dividing both sides by $c^{2}$ and using $\sin A = a/c$, $\cos A = b/c$:
\[
  \left(\frac{a}{c}\right)^{2} + \left(\frac{b}{c}\right)^{2} = 1
  \quad \Longrightarrow \quad
  \sin^{2}(A) + \cos^{2}(A) = 1.
\]
This identity holds for any angle $A$.

\subsection*{Addition and Subtraction Formulas}

The \textbf{cosine addition formula} is
\[
  \cos(A + B) = \cos A \cos B - \sin A \sin B.
\]
The subtraction formula follows by replacing $B$ with $-B$ and using
$\cos(-B) = \cos B$, $\sin(-B) = -\sin B$:
\[
  \cos(A - B) = \cos A \cos B + \sin A \sin B.
\]

\subsection*{Double-Angle Formulas}

\begin{align}
  \sin(2x) &= 2\sin x \cos x, \\
  \cos(2x) &= \cos^{2}x - \sin^{2}x = 2\cos^{2}x - 1 = 1 - 2\sin^{2}x, \\
  \tan(2x) &= \frac{2\tan x}{1 - \tan^{2}x}.
\end{align}

\subsection*{Half-Angle Formulas}

\begin{align}
  \sin\!\left(\frac{x}{2}\right) &= \pm\sqrt{\frac{1 - \cos x}{2}}, \\
  \cos\!\left(\frac{x}{2}\right) &= \pm\sqrt{\frac{1 + \cos x}{2}}, \\
  \tan\!\left(\frac{x}{2}\right)
    &= \pm\sqrt{\frac{1 - \cos x}{1 + \cos x}}
     = \frac{\sin x}{1 + \cos x}
     = \frac{1 - \cos x}{\sin x}.
\end{align}
The correct sign is determined by the quadrant of $x/2$.

\subsection*{Product-to-Sum Formulas}

\begin{align}
  \sin u \sin v &= \tfrac{1}{2}[\cos(u-v) - \cos(u+v)], \\
  \cos u \cos v &= \tfrac{1}{2}[\cos(u-v) + \cos(u+v)], \\
  \sin u \cos v &= \tfrac{1}{2}[\sin(u+v) + \sin(u-v)], \\
  \cos u \sin v &= \tfrac{1}{2}[\sin(u+v) - \sin(u-v)].
\end{align}

\subsection*{Sum-to-Product Formulas}

\begin{align}
  \sin u + \sin v &= 2\sin\!\left(\frac{u+v}{2}\right)\cos\!\left(\frac{u-v}{2}\right), \\
  \sin u - \sin v &= 2\cos\!\left(\frac{u+v}{2}\right)\sin\!\left(\frac{u-v}{2}\right), \\
  \cos u + \cos v &= 2\cos\!\left(\frac{u+v}{2}\right)\cos\!\left(\frac{u-v}{2}\right), \\
  \cos u - \cos v &= -2\sin\!\left(\frac{u+v}{2}\right)\sin\!\left(\frac{u-v}{2}\right).
\end{align}

\subsection*{Law of Sines}

For any triangle with sides $a$, $b$, $c$ opposite angles $A$, $B$, $C$:
\[
  \frac{a}{\sin A} = \frac{b}{\sin B} = \frac{c}{\sin C}.
\]

\begin{itemize}
  \item \href{https://en.wikibooks.org/wiki/A-level_Mathematics/CIE/Pure_Mathematics_2/Trigonometry}{Trigonometry, Wikibooks: Pure Mathematics 2}
\end{itemize}

%-------------------------------------------------------------------------------
\subsection{Learning Objectives}
\label{ssec:lear-fe6e}
%-------------------------------------------------------------------------------

\begin{enumerate}
    \item Use the Pythagorean identity and its corollaries.
    \item Simplify expressions using addition, double-angle, and half-angle formulas.
    \item Apply product-to-sum and sum-to-product formulas.
    \item State and use the Law of Sines.
\end{enumerate}

%-------------------------------------------------------------------------------
\subsection{Assessment Instrument}
\label{ssec:asse-3669}
%-------------------------------------------------------------------------------

Probe understanding of the Pythagorean identity and the addition, double-angle, half-
angle, and product-to-sum formulas through problems that require deriving an unknown trig
value from a given one, verifying identities, and simplifying expressions.

\textbf{Example problems.}
\begin{enumerate}
    \item Given that $\sin A = \frac{3}{5}$ with $A$ in the first quadrant, use the Pythagorean identity to find $\cos A$ and then evaluate $\sin(2A)$ using the double-angle formula.
    \item Use the addition formula $\cos(A - B) = \cos A \cos B + \sin A \sin B$ to find the exact value of $\cos(15^{\circ})$.
    \item Simplify $\frac{1 - \cos(2x)}{\sin(2x)}$ using the half-angle and double-angle formulas, and identify the result as a standard trigonometric function.
\end{enumerate}

%===============================================================================
\section{Section 2.4 - Right Triangle Trigonometry}
\label{sec:sect-974a}
%===============================================================================

This section develops the right-triangle definitions of all six trigonometric functions
and applies them to find exact values at $30^{\circ}$, $45^{\circ}$, and $60^{\circ}$, use
cofunction identities, and solve applied right-triangle problems.

%-------------------------------------------------------------------------------
\subsection{Learning Objectives}
\label{ssec:lear-fe6e}
%-------------------------------------------------------------------------------

\begin{enumerate}
    \item Use right triangles to evaluate trigonometric functions.
    \item Find function values for $30^{\circ}$ ($\pi/6$), $45^{\circ}$ ($\pi/4$), and $60^{\circ}$ ($\pi/3$).
    \item Use cofunctions of complementary angles.
    \item Use the definitions of trigonometric functions of any angle.
    \item Use right triangle trigonometry to solve applied problems.
\end{enumerate}

%-------------------------------------------------------------------------------
\subsection{Assessment Instrument}
\label{ssec:asse-3669}
%-------------------------------------------------------------------------------

Measure competency with right-triangle definitions of the trigonometric functions by
asking students to label triangle sides relative to a given angle, apply the cofunction
identities, and solve applied problems that require setting up and solving a right-
triangle equation.

\textbf{Example problems.}
\begin{enumerate}
    \item In a right triangle with hypotenuse $13$ and one leg $5$, find all six trigonometric functions for the acute angle opposite the side of length $5$.
    \item Use the cofunction identity to evaluate $\sin(62^{\circ})$ in terms of a cosine value, then confirm numerically.
    \item A surveyor stands $50\,\text{m}$ from the base of a building and measures the angle of elevation to the top as $38^{\circ}$. Find the height of the building to the nearest tenth of a metre.
\end{enumerate}

%===============================================================================
\section{Section 3.1 - Graphs of Sine and Cosine Functions}
\label{sec:sect-7809}
%===============================================================================

A \textbf{sinusoidal function} is any function of the form
\[
  f(t) = a\sin(bt + c) + d \quad \text{or} \quad f(t) = a\cos(bt + c) + d,
  \quad a, b > 0.
\]
Three parameters characterise the wave:

\begin{itemize}
  \item \textbf{Amplitude} $a$: the maximum displacement from the midline,
        equal to $|a|$.  The midline (sinusoidal axis) is $y = d$.
  \item \textbf{Period} (wavelength) $T = \frac{2\pi}{b}$: the horizontal
        length of one complete cycle.  The \textbf{frequency} $f = b$ counts
        the number of cycles per $2\pi$ units.
  \item \textbf{Phase shift}: the horizontal displacement, equal to $-c/b$
        (positive leftward when $c > 0$).
\end{itemize}

\textbf{Example.} For $f(t) = 4\sin\!\left(\frac{\pi}{6}t\right) + 40$:
amplitude $= 4$, period $= \frac{2\pi}{\pi/6} = 12$, midline $y = 40$, no
phase shift.

The amplitude can be read from data: if the maximum output is $y_{\max}$ and
the minimum is $y_{\min}$, then
\[
  a = \frac{y_{\max} - y_{\min}}{2}, \qquad d = \frac{y_{\max} + y_{\min}}{2}.
\]
The graphs of $\sin x$ and $\cos x$ have the same shape and differ only by a
phase shift of $\frac{\pi}{2}$: shifting the sine graph $\frac{\pi}{2}$ units
to the left yields the cosine graph.

\begin{itemize}
  \item \href{https://en.wikibooks.org/wiki/Trigonometry/Graphs_of_Sine_and_Cosine_Functions}{Graphs of Sine and Cosine, Wikibooks: Trigonometry}
\end{itemize}

%-------------------------------------------------------------------------------
\subsection{Learning Objectives}
\label{ssec:lear-fe6e}
%-------------------------------------------------------------------------------

\begin{enumerate}
    \item Identify amplitude, period, frequency, phase shift, and midline from the equation of a sinusoidal function.
    \item Sketch the graphs of sine and cosine functions and their transformations.
    \item Write a sinusoidal equation that models a described periodic phenomenon.
\end{enumerate}

%-------------------------------------------------------------------------------
\subsection{Assessment Instrument}
\label{ssec:asse-3669}
%-------------------------------------------------------------------------------

Evaluate the ability to identify and interpret the amplitude, period, frequency, and phase
shift of sinusoidal functions from their equations and from data tables, and to write a
sinusoidal equation that fits a described physical situation.

\textbf{Example problems.}
\begin{enumerate}
    \item For $f(t) = 3\sin\!\left(\frac{\pi}{4}t - \frac{\pi}{2}\right) + 1$, identify the amplitude, period, frequency, phase shift, and midline, then sketch one complete cycle.
    \item A data table shows output values of $g(x)$ at $x = 0, 2, 4, 6, 8$ equal to $5, 0, -5, 0, 5$. Determine the amplitude, period, and an equation of the form $a\cos(bx)$ that fits the data.
    \item Two sinusoidal waves have the same amplitude and frequency but are $90^{\circ}$ out of phase. Write equations for each and verify that their sum is also sinusoidal.
\end{enumerate}

%===============================================================================
\section{Section 3.2 - Graphs of Other Trigonometric Functions}
\label{sec:sect-07b3}
%===============================================================================

The remaining four trigonometric functions are defined as reciprocals or
ratios of sine and cosine; their graphs inherit vertical asymptotes wherever
a denominator is zero.

\subsection*{Tangent}

\[
  \tan x = \frac{\sin x}{\cos x}.
\]
The range of $y = \tan x$ is $(-\infty, \infty)$.  The function is undefined
whenever $\cos x = 0$, i.e., at $x = \frac{\pi}{2} + k\pi$ for any integer
$k$; vertical asymptotes occur at each such value.  The period is $\pi$.

\subsection*{Cotangent}

\[
  \cot x = \frac{\cos x}{\sin x} = \frac{1}{\tan x}.
\]
The range of $y = \cot x$ is also $(-\infty, \infty)$.  Vertical asymptotes
occur at $x = k\pi$ (where $\sin x = 0$).

\subsection*{Secant}

\[
  \sec x = \frac{1}{\cos x}.
\]
Because $\cos x \in [-1, 1]$, the values $\frac{1}{\cos x}$ satisfy
$|\sec x| \geq 1$, giving the range $(-\infty, -1] \cup [1, \infty)$.
Vertical asymptotes coincide with those of $\tan x$.  Secant is an even
function: $\sec(-x) = \sec x$.

\subsection*{Cosecant}

\[
  \csc x = \frac{1}{\sin x}.
\]
The range is the same as that of secant: $(-\infty, -1] \cup [1, \infty)$.
Vertical asymptotes occur at $x = k\pi$.  Cosecant is an odd function:
$\csc(-x) = -\csc x$.

\begin{itemize}
  \item \href{https://en.wikibooks.org/wiki/A-level_Mathematics/CIE/Pure_Mathematics_2/Trigonometry}{Trigonometry, Wikibooks: Pure Mathematics 2}
\end{itemize}

%-------------------------------------------------------------------------------
\subsection{Learning Objectives}
\label{ssec:lear-fe6e}
%-------------------------------------------------------------------------------

\begin{enumerate}
    \item Analyse the graph of $y = \tan x$ and graph its variations.
    \item Analyse the graphs of $y = \sec x$ and $y = \csc x$ and graph their variations.
    \item Analyse the graph of $y = \cot x$ and graph its variations.
\end{enumerate}

%-------------------------------------------------------------------------------
\subsection{Assessment Instrument}
\label{ssec:asse-3669}
%-------------------------------------------------------------------------------

Assess understanding of the graphs of tangent, cotangent, secant, and cosecant by
requiring students to identify vertical asymptotes, domains, ranges, and periods, and to
sketch graphs over a specified interval.

\textbf{Example problems.}
\begin{enumerate}
    \item State the domain, range, and period of $y = \tan x$ and identify all vertical asymptotes in $[-\pi, \pi]$.
    \item Sketch $y = \sec x$ and $y = \cos x$ on the same axes over $[-2\pi, 2\pi]$. Mark the vertical asymptotes of $\sec x$ and explain why they occur at those positions.
    \item Explain why $y = \csc x$ has the same range as $y = \sec x$ but different asymptotes, and give the general form of the asymptotes of each.
\end{enumerate}

%===============================================================================
\section{Section 3.3 - Inverse Functions}
\label{sec:sect-f100}
%===============================================================================

The function $g$ is the \textbf{inverse} of the one-to-one function $f$ if
and only if
\[
  g(f(x)) = x \quad \text{for all } x \text{ in the domain of } f, \qquad
  f(g(x)) = x \quad \text{for all } x \text{ in the domain of } g.
\]
The inverse is denoted $f^{-1}$.  Geometrically, the graph of $f^{-1}$ is the
reflection of the graph of $f$ across the line $y = x$.

A function has an inverse if and only if it is \textbf{one-to-one}: the
horizontal line test determines this visually.

\subsection*{Finding an Inverse Algebraically}

Write $x = f(f^{-1}(x))$ and solve for $f^{-1}(x)$.

\textbf{Example 1.} Let $f(x) = 2x + 1$.  Then
\begin{align}
  x &= 2f^{-1}(x) + 1 \\
  f^{-1}(x) &= \frac{x - 1}{2}.
\end{align}

\textbf{Example 2.} Let $f(x) = \sqrt{x - 1}$ (domain $x \geq 1$).  Then
\begin{align}
  x &= \sqrt{f^{-1}(x) - 1} \\
  x^{2} &= f^{-1}(x) - 1 \\
  f^{-1}(x) &= x^{2} + 1, \quad x \geq 0.
\end{align}
The domain of $f^{-1}$ equals the range of $f$, so the restriction $x \geq 0$
is necessary.

\begin{itemize}
  \item \href{https://en.wikibooks.org/wiki/Algebra/Functions}{Functions, Wikibooks: Algebra}
  \item \href{https://en.wikibooks.org/wiki/VCE_Mathematical_Methods/Inverse_functions}{Inverse Functions, Wikibooks: VCE Mathematical Methods}
\end{itemize}

%-------------------------------------------------------------------------------
\subsection{Learning Objectives}
\label{ssec:lear-fe6e}
%-------------------------------------------------------------------------------

\begin{enumerate}
    \item Verify inverse functions.
    \item Determine the domain and range of an inverse function, and restrict the domain of a function to make it one-to-one.
    \item Find or evaluate the inverse of a function.
    \item Use the graph of a one-to-one function to graph its inverse function on the same axes.
\end{enumerate}

%-------------------------------------------------------------------------------
\subsection{Assessment Instrument}
\label{ssec:asse-3669}
%-------------------------------------------------------------------------------

Test the ability to find inverse functions algebraically, verify inverses using the
composition test, determine the domain and range of an inverse, and apply the horizontal
line test to decide whether an inverse function exists.

\textbf{Example problems.}
\begin{enumerate}
    \item Find the inverse of $f(x) = \frac{x - 3}{2x + 1}$ and verify that $f(f^{-1}(x)) = x$.
    \item The function $f(x) = x^{2} - 4$ is not one-to-one on $\mathbb{R}$. Restrict its domain so that an inverse exists, then find and state the domain of $f^{-1}$.
    \item Given $f(x) = \sqrt{x - 1}$ (with the natural domain), find $f^{-1}(x)$, state its domain, and sketch both $f$ and $f^{-1}$ on the same axes, labelling their relationship to the line $y = x$.
\end{enumerate}

%===============================================================================
\section{Section 3.4 - Inverse Trigonometric Functions}
\label{sec:sect-6b22}
%===============================================================================

Because no trigonometric function is one-to-one on its full domain, each must
be restricted to a \textbf{principal branch} before an inverse can be defined.
The resulting \textbf{inverse trigonometric functions} return a unique
\textbf{principal value} for each input.

\subsection*{Notation}

The preferred notation uses the arc-prefix: $\arcsin x$, $\arccos x$,
$\arctan x$, and so on.  The notation $\sin^{-1}(x)$, introduced by John
Herschel in 1813, is also common but can be confused with the reciprocal
$(\sin x)^{-1} = \csc x$; the ISO 80000-2 standard (2009) specifies the
arc-prefix notation exclusively.

\subsection*{Principal Domains and Ranges}

\[
  \begin{array}{|l|c|c|}
    \hline
    \textbf{Function} & \textbf{Domain of } x & \textbf{Principal range (rad)} \\
    \hline
    y = \arcsin x  & [-1,\,1] & [-\pi/2,\;\pi/2] \\
    y = \arccos x  & [-1,\,1] & [0,\;\pi] \\
    y = \arctan x  & \mathbb{R} & (-\pi/2,\;\pi/2) \\
    y = \operatorname{arccot} x & \mathbb{R} & (0,\;\pi) \\
    y = \operatorname{arcsec} x & |x| \geq 1 & [0,\pi]\setminus\{\pi/2\} \\
    y = \operatorname{arccsc} x & |x| \geq 1 & [-\pi/2,\pi/2]\setminus\{0\} \\
    \hline
  \end{array}
\]

\subsection*{Compositions of Trigonometric and Inverse Trigonometric
Functions}

A right-triangle diagram with one leg $x$ and hypotenuse $1$ (or
hypotenuse $\sqrt{1+x^{2}}$ for arctangent) gives:
\[
  \cos(\arcsin x) = \sqrt{1 - x^{2}}, \quad
  \sin(\arccos x) = \sqrt{1 - x^{2}}, \quad
  \tan(\arcsin x) = \frac{x}{\sqrt{1 - x^{2}}}.
\]

\subsection*{Solving Equations Using Inverse Functions}

The equation $\cos\theta = x$ has the general solution
\[
  \theta = \pm\arccos x + 2\pi k, \quad k \in \mathbb{Z}.
\]
The $\pm$ symbol means that both $+\arccos x$ and $-\arccos x$ (plus any
integer multiple of $2\pi$) are solutions; additional context is required to
determine which branch applies in a given problem.

\begin{itemize}
  \item \href{https://en.wikibooks.org/wiki/Trigonometry/Inverse_Trigonometric_Functions}{Inverse Trig Functions, Wikibooks: Trigonometry}
\end{itemize}

%-------------------------------------------------------------------------------
\subsection{Learning Objectives}
\label{ssec:lear-fe6e}
%-------------------------------------------------------------------------------

\begin{enumerate}
    \item Understand the notation and principal-value restrictions of the six inverse trigonometric functions.
    \item Evaluate expressions of the form $\arcsin x$, $\arccos x$, and $\arctan x$.
    \item Simplify compositions of trigonometric and inverse trigonometric functions.
    \item Solve elementary trigonometric equations using inverse functions.
\end{enumerate}

%-------------------------------------------------------------------------------
\subsection{Assessment Instrument}
\label{ssec:asse-3669}
%-------------------------------------------------------------------------------

Evaluate mastery of inverse trigonometric notation, principal-value restrictions, and the
composition identities by asking students to evaluate expressions, solve elementary
trigonometric equations, and simplify compositions such as $\sin(\arccos x)$.

\textbf{Example problems.}
\begin{enumerate}
    \item Evaluate $\arcsin\!\left(-\frac{\sqrt{2}}{2}\right)$, $\arccos(0)$, and $\arctan(1)$, expressing each answer in radians and degrees.
    \item Simplify $\cos(\arcsin x)$ for $x \in [-1, 1]$ by drawing an appropriate right triangle, and verify your result using the Pythagorean theorem.
    \item Solve $\cos\theta = -\frac{1}{2}$ for all $\theta \in [0, 2\pi)$, expressing the solutions using the inverse cosine function and explaining the role of the $\pm$ symbol.
\end{enumerate}

%===============================================================================
\section{Section 4.1 - Solving Trigonometric Equations with Identities}
\label{sec:sect-fb9f}
%===============================================================================

The three \textbf{Pythagorean identities} follow from $\sin^{2}A + \cos^{2}A
= 1$ by dividing through by $\cos^{2}A$ or $\sin^{2}A$:
\[
  \sin^{2}A + \cos^{2}A = 1, \qquad
  1 + \tan^{2}A = \sec^{2}A, \qquad
  1 + \cot^{2}A = \csc^{2}A.
\]
These identities, together with the reciprocal and quotient identities, allow
complex trigonometric expressions to be rewritten in simpler forms before
solving.

\textbf{Strategy.} To verify an identity, transform one side (usually the more
complex one) into the other using known identities, one step at a time, without
modifying both sides simultaneously.  Common techniques include replacing
$\sec x$ with $\frac{1}{\cos x}$, factoring a Pythagorean expression, and
multiplying by a conjugate.

\textbf{Example.} Verify $\frac{1 - \cos^{2}x}{\cos x} = \sin x \tan x$.

Starting from the left-hand side and applying $\sin^{2}x = 1 - \cos^{2}x$:
\[
  \frac{1 - \cos^{2}x}{\cos x}
  = \frac{\sin^{2}x}{\cos x}
  = \sin x \cdot \frac{\sin x}{\cos x}
  = \sin x \tan x. \qquad \checkmark
\]

\begin{itemize}
  \item \href{https://en.wikibooks.org/wiki/A-level_Mathematics/CIE/Pure_Mathematics_2/Trigonometry}{Trigonometry, Wikibooks: Pure Mathematics 2}
\end{itemize}

%-------------------------------------------------------------------------------
\subsection{Learning Objectives}
\label{ssec:lear-fe6e}
%-------------------------------------------------------------------------------

\begin{enumerate}
    \item Verify trigonometric identities.
    \item Simplify trigonometric expressions using the Pythagorean, reciprocal, and quotient identities.
    \item Use identities as a tool to solve trigonometric equations.
\end{enumerate}

%-------------------------------------------------------------------------------
\subsection{Assessment Instrument}
\label{ssec:asse-3669}
%-------------------------------------------------------------------------------

Assess the ability to use the Pythagorean, reciprocal, and quotient identities to verify
equations and to simplify expressions before solving them, emphasising multi-step
reasoning and justification of each algebraic transformation.

\textbf{Example problems.}
\begin{enumerate}
    \item Verify the identity $\frac{\sin^{2}\theta}{1 - \cos\theta} = 1 + \cos\theta$ by transforming the left-hand side.
    \item Use the Pythagorean identity to simplify $\frac{1 - \sin^{2}x}{\cos x} + \frac{1 - \cos^{2}x}{\sin x}$ to a single trigonometric expression.
    \item For a right triangle with hypotenuse $c$ and legs $a$ and $b$, derive the identity $\sin^{2}A + \cos^{2}A = 1$ from the Pythagorean theorem, where $A$ is the angle opposite side $a$.
\end{enumerate}

%===============================================================================
\section{Section 4.2 - Sum and Difference Identities}
\label{sec:sect-8465}
%===============================================================================

This section establishes and applies the sum and difference formulas for cosine, sine, and
tangent, including their use with cofunction relationships and in verifying identities.

%-------------------------------------------------------------------------------
\subsection{Learning Objectives}
\label{ssec:lear-fe6e}
%-------------------------------------------------------------------------------

\begin{enumerate}
    \item Use sum and difference formulas for cosine.
    \item Use sum and difference formulas for sine.
    \item Use sum and difference formulas for tangent.
    \item Use sum and difference formulas for cofunctions.
    \item Use sum and difference formulas to verify identities.
\end{enumerate}

%-------------------------------------------------------------------------------
\subsection{Assessment Instrument}
\label{ssec:asse-3669}
%-------------------------------------------------------------------------------

Probe fluency with the sum and difference formulas for sine, cosine, and tangent by
requiring exact-value calculations at non-standard angles, cofunction applications, and
identity verifications.

\textbf{Example problems.}
\begin{enumerate}
    \item Use the cosine difference formula to find the exact value of $\cos(75^{\circ})$.
    \item Use the sine addition formula to verify the cofunction identity $\sin\!\left(\frac{\pi}{2} - x\right) = \cos x$.
    \item Use the tangent addition formula to show that $\tan\!\left(x + \frac{\pi}{4}\right) = \frac{1 + \tan x}{1 - \tan x}$.
\end{enumerate}

%===============================================================================
\section{Section 4.3 - Double-Angle, Half-Angle, and Reduction Formulas}
\label{sec:sect-b254}
%===============================================================================

This section derives and applies the double-angle formulas, reduction formulas, and half-
angle formulas to find exact values and to verify trigonometric identities.

%-------------------------------------------------------------------------------
\subsection{Learning Objectives}
\label{ssec:lear-fe6e}
%-------------------------------------------------------------------------------

\begin{enumerate}
    \item Use double-angle formulas to find exact values.
    \item Use double-angle formulas to verify identities.
    \item Use reduction formulas to simplify an expression.
    \item Use half-angle formulas to find exact values.
\end{enumerate}

%-------------------------------------------------------------------------------
\subsection{Assessment Instrument}
\label{ssec:asse-3669}
%-------------------------------------------------------------------------------

Test the ability to apply double-angle and half-angle formulas to find exact values,
simplify expressions, and verify identities, including choosing the correct sign for the
half-angle formula based on the quadrant of the argument.

\textbf{Example problems.}
\begin{enumerate}
    \item Given $\cos\theta = \frac{5}{13}$ with $\theta$ in the first quadrant, find $\sin(2\theta)$, $\cos(2\theta)$, and $\cos\!\left(\frac{\theta}{2}\right)$.
    \item Use the double-angle formula to verify the identity $\frac{\sin(2x)}{1 + \cos(2x)} = \tan x$.
    \item Find the exact value of $\sin\!\left(\frac{\pi}{8}\right)$ using the half-angle formula, specifying why you choose the positive square root.
\end{enumerate}

%===============================================================================
\section{Section 4.4 - Sum-to-Product and Product-to-Sum Formulas}
\label{sec:sect-f015}
%===============================================================================

This section presents the product-to-sum and sum-to-product conversion formulas and
develops the ability to rewrite trigonometric expressions in either product or sum form.

%-------------------------------------------------------------------------------
\subsection{Learning Objectives}
\label{ssec:lear-fe6e}
%-------------------------------------------------------------------------------

\begin{enumerate}
    \item Express products as sums.
    \item Express sums as products.
\end{enumerate}

%-------------------------------------------------------------------------------
\subsection{Assessment Instrument}
\label{ssec:asse-3669}
%-------------------------------------------------------------------------------

Assess the ability to convert between product and sum forms of trigonometric expressions
by applying the relevant identities in both directions, and to use these conversions to
simplify or evaluate expressions.

\textbf{Example problems.}
\begin{enumerate}
    \item Express $\sin(5x)\cos(3x)$ as a sum of trigonometric functions.
    \item Express $\sin(4x) + \sin(2x)$ as a product of trigonometric functions.
    \item Use a product-to-sum identity to evaluate $2\sin(75^{\circ})\sin(15^{\circ})$ exactly.
\end{enumerate}

%===============================================================================
\section{Section 4.5 - Solving Trigonometric Equations}
\label{sec:sect-8eb2}
%===============================================================================

This section addresses the full range of trigonometric equation types: linear, quadratic-
in-form, multi-angle, and those requiring an identity substitution, solved both
analytically and with a calculator.

%-------------------------------------------------------------------------------
\subsection{Learning Objectives}
\label{ssec:lear-fe6e}
%-------------------------------------------------------------------------------

\begin{enumerate}
    \item Solve linear trigonometric equations in sine and cosine.
    \item Solve equations involving a single trigonometric function.
    \item Solve trigonometric equations using a calculator.
    \item Solve trigonometric equations that are quadratic in form.
    \item Solve trigonometric equations using fundamental identities.
    \item Solve trigonometric equations with multiple angles.
    \item Solve right triangle problems.
\end{enumerate}

%-------------------------------------------------------------------------------
\subsection{Assessment Instrument}
\label{ssec:asse-3669}
%-------------------------------------------------------------------------------

Evaluate the ability to solve trigonometric equations over a given interval and over all
real numbers, including linear, quadratic-in-form, and multi-angle equations, and to use a
calculator where exact solutions are not available.

\textbf{Example problems.}
\begin{enumerate}
    \item Solve $2\cos\theta - \sqrt{3} = 0$ for all $\theta \in [0, 2\pi)$.
    \item Solve the quadratic equation $2\sin^{2}x + \sin x - 1 = 0$ for $x \in [0^{\circ}, 360^{\circ})$.
    \item Solve $\sin(2x) = \cos x$ for $x \in [0, 2\pi)$ by using an appropriate identity to reduce the equation to a single trigonometric function.
\end{enumerate}

%===============================================================================
\section{Section 5.1 - Non-Right Triangles: Law of Sines}
\label{sec:sect-afc1}
%===============================================================================

This section introduces the Law of Sines for solving oblique triangles in AAS, ASA, and
SSA (ambiguous) configurations, and derives the formula for triangle area using the sine
function.

%-------------------------------------------------------------------------------
\subsection{Learning Objectives}
\label{ssec:lear-fe6e}
%-------------------------------------------------------------------------------

\begin{enumerate}
    \item Use the Law of Sines to solve oblique triangles.
    \item Find the area of an oblique triangle using the sine function.
    \item Solve applied problems using the Law of Sines.
\end{enumerate}

%-------------------------------------------------------------------------------
\subsection{Assessment Instrument}
\label{ssec:asse-3669}
%-------------------------------------------------------------------------------

Assess the ability to determine when the Law of Sines applies, to solve for unknown sides
and angles in oblique triangles (including the ambiguous SSA case), and to compute
triangle area from two sides and the included angle.

\textbf{Example problems.}
\begin{enumerate}
    \item In triangle $ABC$, $A = 35^{\circ}$, $B = 85^{\circ}$, and $a = 14$. Find sides $b$ and $c$.
    \item Two angles and the side opposite one of them are given: $A = 40^{\circ}$, $a = 10$, $b = 14$. Determine how many triangles are possible and solve each.
    \item Find the area of triangle $ABC$ given $a = 7$, $b = 9$, and $C = 120^{\circ}$.
\end{enumerate}

%===============================================================================
\section{Section 5.2 - Non-Right Triangles: Law of Cosines}
\label{sec:sect-800a}
%===============================================================================

This section presents the Law of Cosines for the SAS and SSS cases, applies it to
surveying and navigation problems, and introduces Heron's formula as an alternative area
method.

%-------------------------------------------------------------------------------
\subsection{Learning Objectives}
\label{ssec:lear-fe6e}
%-------------------------------------------------------------------------------

\begin{enumerate}
    \item Use the Law of Cosines to solve oblique triangles.
    \item Solve applied problems using the Law of Cosines.
    \item Use Heron's formula to find the area of a triangle.
\end{enumerate}

%-------------------------------------------------------------------------------
\subsection{Assessment Instrument}
\label{ssec:asse-3669}
%-------------------------------------------------------------------------------

Test the ability to select and apply the Law of Cosines for the SAS and SSS cases, to
solve applied navigation or surveying problems, and to use Heron's formula as an
alternative area computation when all three sides are known.

\textbf{Example problems.}
\begin{enumerate}
    \item In triangle $PQR$, $p = 8$, $q = 11$, and $R = 60^{\circ}$. Find side $r$ and the remaining angles.
    \item Given a triangle with sides $a = 5$, $b = 7$, and $c = 10$, find all three angles using the Law of Cosines.
    \item Use Heron's formula to find the area of the triangle in the previous problem, and verify using the formula $K = \frac{1}{2}ab\sin C$.
\end{enumerate}

%===============================================================================
\section{Section 5.3 - Polar Coordinates}
\label{sec:sect-1544}
%===============================================================================

This section introduces polar coordinates, establishes the conversion equations between
polar and rectangular forms, and practises transforming equations and identifying standard
curves.

%-------------------------------------------------------------------------------
\subsection{Learning Objectives}
\label{ssec:lear-fe6e}
%-------------------------------------------------------------------------------

\begin{enumerate}
    \item Plot points using polar coordinates.
    \item Convert from polar coordinates to rectangular coordinates.
    \item Convert from rectangular coordinates to polar coordinates.
    \item Transform equations between polar and rectangular forms.
    \item Identify and graph polar equations by converting to rectangular equations.
\end{enumerate}

%-------------------------------------------------------------------------------
\subsection{Assessment Instrument}
\label{ssec:asse-3669}
%-------------------------------------------------------------------------------

Evaluate fluency with the polar coordinate system by asking students to plot points,
convert between polar and rectangular forms of points and equations, and identify standard
curves from their polar equations.

\textbf{Example problems.}
\begin{enumerate}
    \item Convert the rectangular point $(-3, 3)$ to polar coordinates $(r, \theta)$ with $r > 0$ and $0 \leq \theta < 2\pi$.
    \item Convert the polar equation $r = 4\sin\theta$ to rectangular form and identify the resulting curve.
    \item Convert the rectangular equation $x^{2} + y^{2} = 9$ to polar form and sketch the curve.
\end{enumerate}

%===============================================================================
\section{Section 5.4 - Polar Coordinates: Graphs}
\label{sec:sect-5223}
%===============================================================================

This section develops the ability to test polar equations for symmetry and to graph common
polar curves by plotting key points.

%-------------------------------------------------------------------------------
\subsection{Learning Objectives}
\label{ssec:lear-fe6e}
%-------------------------------------------------------------------------------

\begin{enumerate}
    \item Test polar equations for symmetry.
    \item Graph polar equations by plotting points.
\end{enumerate}

%-------------------------------------------------------------------------------
\subsection{Assessment Instrument}
\label{ssec:asse-3669}
%-------------------------------------------------------------------------------

Assess the ability to identify symmetry in polar equations and to graph common polar
curves (circles, limaçons, roses, lemniscates) by plotting key points and using the
symmetry tests.

\textbf{Example problems.}
\begin{enumerate}
    \item Test the equation $r = 3 + 2\cos\theta$ for symmetry with respect to the polar axis and sketch the curve.
    \item Sketch the rose curve $r = 2\cos(3\theta)$ by finding the angles where $r = 0$ and where $r$ is maximised.
    \item Determine the symmetry of $r^{2} = 9\cos(2\theta)$ and sketch the lemniscate.
\end{enumerate}

%===============================================================================
\section{Section 5.5 - The Polar Form of Complex Numbers}
\label{sec:sect-9c90}
%===============================================================================

This section connects the polar coordinate system to complex numbers, covering conversion
to polar form, multiplication and division of complex numbers in polar form, and De
Moivre's theorem for powers and roots.

%-------------------------------------------------------------------------------
\subsection{Learning Objectives}
\label{ssec:lear-fe6e}
%-------------------------------------------------------------------------------

\begin{enumerate}
    \item Plot complex numbers in the complex plane.
    \item Find the absolute value of a complex number.
    \item Write complex numbers in polar form.
    \item Convert a complex number from polar to rectangular form.
    \item Find products of complex numbers in polar form.
    \item Find quotients of complex numbers in polar form.
    \item Find powers of complex numbers in polar form.
    \item Find roots of complex numbers in polar form.
\end{enumerate}

%-------------------------------------------------------------------------------
\subsection{Assessment Instrument}
\label{ssec:asse-3669}
%-------------------------------------------------------------------------------

Test the ability to convert complex numbers between rectangular and polar form, to
multiply and divide using moduli and arguments, and to apply De Moivre's theorem to
compute powers and roots of complex numbers.

\textbf{Example problems.}
\begin{enumerate}
    \item Write $z = -1 + \sqrt{3}\,i$ in polar form $r(\cos\theta + i\sin\theta)$ and find $|z|$ and $\arg z$.
    \item Let $z_{1} = 2(\cos 30^{\circ} + i\sin 30^{\circ})$ and $z_{2} = 3(\cos 90^{\circ} + i\sin 90^{\circ})$. Find $z_{1}z_{2}$ and $z_{1}/z_{2}$ in polar form.
    \item Use De Moivre's theorem to find all cube roots of $8i$ and express them in rectangular form.
\end{enumerate}

%===============================================================================
\section{Section 6.1 - Exponential Functions and Their Graphs}
\label{sec:sect-6796}
%===============================================================================

This section introduces exponential functions, their evaluation and graphical properties,
the compound-interest formula, and the natural base $e$.

%-------------------------------------------------------------------------------
\subsection{Learning Objectives}
\label{ssec:lear-fe6e}
%-------------------------------------------------------------------------------

\begin{enumerate}
    \item Evaluate exponential functions.
    \item Find the equation of an exponential function.
    \item Use compound interest formulas.
    \item Evaluate exponential functions with base $e$.
    \item Graph exponential functions.
    \item Graph exponential functions using transformations.
\end{enumerate}

%-------------------------------------------------------------------------------
\subsection{Assessment Instrument}
\label{ssec:asse-3669}
%-------------------------------------------------------------------------------

Assess understanding of exponential functions through problems requiring evaluation, graph
transformations, and application of compound interest formulas, including the continuously
compounding case with base $e$.

\textbf{Example problems.}
\begin{enumerate}
    \item Evaluate $f(x) = 3^{x}$ at $x = -2, 0, 1, 2$ and sketch the graph, labelling the $y$-intercept and the horizontal asymptote.
    \item An investment of \$1{,}000 is compounded quarterly at an annual rate of $6\%$. Write the compound-interest formula and find the balance after $5$ years.
    \item Describe the transformation from $f(x) = e^{x}$ to $g(x) = 2e^{x-3} + 1$ and sketch both curves on the same axes.
\end{enumerate}

%===============================================================================
\section{Section 6.2 - Logarithmic Functions and Their Graphs}
\label{sec:sect-84b2}
%===============================================================================

This section covers the conversion between exponential and logarithmic form, evaluation of
common and natural logarithms, and the domain, range, and graph of logarithmic functions.

%-------------------------------------------------------------------------------
\subsection{Learning Objectives}
\label{ssec:lear-fe6e}
%-------------------------------------------------------------------------------

\begin{enumerate}
    \item Convert from logarithmic to exponential form.
    \item Convert from exponential to logarithmic form.
    \item Evaluate logarithms.
    \item Use common logarithms.
    \item Use natural logarithms.
    \item Identify the domain and range of a logarithmic function.
    \item Graph logarithmic functions.
\end{enumerate}

%-------------------------------------------------------------------------------
\subsection{Assessment Instrument}
\label{ssec:asse-3669}
%-------------------------------------------------------------------------------

Evaluate the ability to convert between exponential and logarithmic form, to evaluate
common and natural logarithms, and to identify the domain, range, and graph of a
logarithmic function including its transformations.

\textbf{Example problems.}
\begin{enumerate}
    \item Rewrite $2^{5} = 32$ in logarithmic form and evaluate $\log_{3}(81)$ and $\ln(e^{4})$.
    \item Find the domain of $f(x) = \log_{2}(3x - 6)$ and sketch the graph, indicating the vertical asymptote and the $x$-intercept.
    \item Describe the transformation from $f(x) = \log x$ to $g(x) = -\log(x + 2) + 3$ and state the domain and range of $g$.
\end{enumerate}

%===============================================================================
\section{Section 6.3 - Logarithmic Properties}
\label{sec:sect-9905}
%===============================================================================

This section develops the product, quotient, and power rules for logarithms and applies
them to expand, condense, and re-base logarithmic expressions.

%-------------------------------------------------------------------------------
\subsection{Learning Objectives}
\label{ssec:lear-fe6e}
%-------------------------------------------------------------------------------

\begin{enumerate}
    \item Use the product rule for logarithms.
    \item Use the quotient rule for logarithms.
    \item Use the power rule for logarithms.
    \item Expand logarithmic expressions.
    \item Condense logarithmic expressions.
    \item Use the change-of-base formula for logarithms.
\end{enumerate}

%-------------------------------------------------------------------------------
\subsection{Assessment Instrument}
\label{ssec:asse-3669}
%-------------------------------------------------------------------------------

Test mastery of the product, quotient, and power rules for logarithms by asking students
to expand and condense logarithmic expressions and to apply the change-of-base formula to
evaluate non-standard logarithms.

\textbf{Example problems.}
\begin{enumerate}
    \item Expand $\log_{2}\!\left(\frac{x^{3}\sqrt{y}}{z^{2}}\right)$ using the product, quotient, and power rules.
    \item Condense $3\ln x - \frac{1}{2}\ln y + 2\ln z$ into a single logarithm.
    \item Use the change-of-base formula to evaluate $\log_{7}(200)$ to four decimal places.
\end{enumerate}

%===============================================================================
\section{Section 6.4 - Exponential and Logarithmic Equations}
\label{sec:sect-64ce}
%===============================================================================

This section presents strategies for solving exponential and logarithmic equations,
including the like-base method, the one-to-one property, and applications to growth,
decay, and finance.

%-------------------------------------------------------------------------------
\subsection{Learning Objectives}
\label{ssec:lear-fe6e}
%-------------------------------------------------------------------------------

\begin{enumerate}
    \item Use like bases to solve exponential equations.
    \item Use logarithms to solve exponential equations.
    \item Use the definition of a logarithm to solve logarithmic equations.
    \item Use the one-to-one property of logarithms to solve logarithmic equations.
    \item Solve applied problems involving exponential and logarithmic equations.
\end{enumerate}

%-------------------------------------------------------------------------------
\subsection{Assessment Instrument}
\label{ssec:asse-3669}
%-------------------------------------------------------------------------------

Assess the ability to solve exponential and logarithmic equations using like-base
strategies, logarithm application, and the one-to-one property, including applied problems
such as exponential growth, decay, and finance.

\textbf{Example problems.}
\begin{enumerate}
    \item Solve $4^{x+1} = 64$ using like bases, and verify by substitution.
    \item Solve $\ln(2x + 3) = 5$ for $x$ and check that the solution is in the domain of the original equation.
    \item A radioactive substance decays according to $N(t) = N_{0}e^{-0.03t}$. Find the half-life of the substance and determine how long it takes for the sample to decay to $10\%$ of its initial amount.
\end{enumerate}